\begin{equation}\tag{6.0.3}
\bar\alpha_S:W_S\longrightarrow W_S
\end{equation}

\noindent
表示 $\alpha_S$ 诱导的 $\tau^{-1}$ 线性运算;

\begin{equation}\tag{6.0.4}
R_S(i)
=\text{$R_S$ 中关于 $X_0$ 次数为 $i$ 的元素组成的子模},
\end{equation}

且

\begin{equation}\tag{6.0.5}
W_S(i)=R_S(i)\otimes_{\mathcal O}\mathbb F_q.
\end{equation}

\subsubsection*{命题 6.0.6.}
有

\[
\alpha_S(R_S)\subset R_S(1)+pR_S.
\]

\noindent\emph{证明.}
有分解 (4.6.14)

\begin{equation}\tag{6.0.7}
\alpha_S=\frac1p\,\tau^{-1}\circ\psi_p\circ H(X).
\end{equation}

已知 (见 (4.6.15))

\begin{equation}\tag{6.0.8}
H(X)\,R_S\subset R_S,
\end{equation}

且显然

\begin{equation}\tag{6.0.9}
\tau^{-1}\bigl(R_S(1)+pR_S\bigr)=R_S(1)+pR_S.
\end{equation}

因此只需验证

\begin{equation}\tag{6.0.10}
\frac1p\,\psi_p(R_S)\subset R_S(1)+pR_S,
\end{equation}

这由 (4.6.19) 显然成立.

\subsubsection*{推论 6.0.11.}
有
\[
\bar\alpha_S(W_S)\subset W_S(1).
\]

\subsubsection*{推论 6.0.12.}
有同余式

\[
\Delta_S(t)\equiv
\det\!\left(
1-t(\bar\alpha_S)^{\operatorname{ord}(q)}
\mathrel{\Big|}W_S(1)\right)
\pmod{\pi\mathcal O[[t]]}.
\]
